% 117-04(117쪽) — 삼차함수 y=f(x) 의 그래프(x축과 -3 에서 만나고 원점에서 접한다)와 -3~0 사이의 색칠한 부분. 스캔 화질이 낮아 TikZ 로 다시 그림.
% 원문에 있는 것만: 곡선 · 색칠한 부분 · 눈금 -3 · 라벨 O, x, y, y=f(x). y축 눈금과 넓이 값은 원문에 없으므로 넣지 않는다.
\begin{tikzpicture}[x=0.70cm, y=0.38cm, line width=0.5pt]
  % 색칠한 부분(-3 ≤ x ≤ 0)
  \fill[gray!18] plot[smooth, tension=0.6] coordinates {
    (-3,0) (-2.90,0.84) (-2.75,1.89) (-2.60,2.70) (-2.40,3.46) (-2.20,3.87)
    (-2.00,4.00) (-1.80,3.89) (-1.60,3.58) (-1.40,3.14) (-1.20,2.59) (-1.00,2.00)
    (-0.80,1.41) (-0.60,0.86) (-0.40,0.42) (-0.20,0.11) (0,0) } -- cycle;
  \draw[->, >=stealth, line width=0.7pt] (-3.80,0) -- (1.60,0) node[below=1pt, font=\small] {$x$};
  \draw[->, >=stealth, line width=0.7pt] (0,-2.10) -- (0,6.70) node[left=1pt, font=\small] {$y$};
  % 곡선
  \draw[line width=0.6pt] plot[smooth, tension=0.6] coordinates {
    (-3.12,-1.17) (-3.05,-0.47) (-3,0) (-2.90,0.84) (-2.75,1.89) (-2.60,2.70)
    (-2.40,3.46) (-2.20,3.87) (-2.00,4.00) (-1.80,3.89) (-1.60,3.58) (-1.40,3.14)
    (-1.20,2.59) (-1.00,2.00) (-0.80,1.41) (-0.60,0.86) (-0.40,0.42) (-0.20,0.11)
    (0,0) (0.20,0.13) (0.40,0.54) (0.60,1.30) (0.80,2.43) (1.00,4.00) (1.15,5.49) };
  \node[below=2pt, font=\small] at (-3.32,0) {$-3$};
  \node[below left=0pt and -2pt, font=\small] at (0,0) {$\pt{O}$};
  \node[anchor=west, font=\small] at (0.14,6.10) {$y=f(x)$};
\end{tikzpicture}
